% GENERATED FILE. Edit canonical lessons, then run npm run generate:books.
% canonical-source-hash: fnv1a64:a13a8f232b164ced
% canonical-lessons: FA-C45-khane, FA-C45-otagh, FA-C45-panjere, FA-C45-sandali, FA-C45-miz

\chapter{House, Room, Window, Chair, Table}
\label{ch:fa-house-room-window-chair-table}
\hlchaptermodality{45}

\begin{chapteropening}
\textbf{By the end of this chapter:} \emph{I can say a house, home, a room, a window, a chair and a table.}

Complete the last lesson of chapter 45: I can say a house, home, a room, a window, a chair and a table.
\end{chapteropening}

\section[\texorpdfstring{\fa{خانه} (khâne)}{خانه (khâne)}]{\fa{خانه} (khâne) — a house, home}

\label{lesson:FA-C45-khane}



\begin{quote}
Before the new one: say the Persian for cheese, then the Persian for milk; a lion.
\end{quote}

\subsection*{You'll want to know: \fa{خانه}}
\textbf{\fa{خانه}} — \emph{khâne} — \textquotedblleft{}a house, home\textquotedblright{}.

\begin{cousinweb}
\emph{khâne} is a native Persian word, inherited rather than borrowed; it is usually traced to Middle Persian \emph{xānag}, \textquotedblleft{}house\textquotedblright{}.
\end{cousinweb}

\begin{grammarlens}[title={What you've built}]
One more word for this chapter.
\end{grammarlens}

\subsection*{Your turn}
\begin{itemize}
\raggedright
  \item \emph{Say it:} \emph{khâne}
  \item \emph{Say it:} \emph{khâne}, once more
  \item \emph{Say it:} say \emph{khâne}
  \item \emph{Recall:} say the Persian for cheese, then the Persian for milk; a lion, then say \emph{khâne} again
\end{itemize}

\begin{tcolorbox}[breakable,colback=teal!4,colframe=teal!35!black,title={Before you move on}]
What does \fa{خانه} mean? (\textquotedblleft{}A house, home\textquotedblright{}.) Say it once more.
\end{tcolorbox}
\section[\texorpdfstring{\fa{اتاق} (otâgh)}{اتاق (otâgh)}]{\fa{اتاق} (otâgh) — a room}

\label{lesson:FA-C45-otagh}



\begin{quote}
Before the new one: say the Persian for milk; a lion, then the Persian for a house.
\end{quote}

\subsection*{You'll want to know: \fa{اتاق}}
\textbf{\fa{اتاق}} — \emph{otâgh} — \textquotedblleft{}a room\textquotedblright{}.

\begin{cousinweb}
\emph{otâgh} is a loan, usually traced to Turkic \emph{otaq}, \textquotedblleft{}tent, room\textquotedblright{}.
\end{cousinweb}

\begin{grammarlens}[title={What you've built}]
One more word for this chapter.
\end{grammarlens}

\subsection*{Your turn}
\begin{itemize}
\raggedright
  \item \emph{Say it:} \emph{otâgh}
  \item \emph{Say it:} \emph{otâgh}, once more
  \item \emph{Say it:} say \emph{otâgh}
  \item \emph{Recall:} say the Persian for milk; a lion, then the Persian for a house, then say \emph{otâgh} again
\end{itemize}

\begin{tcolorbox}[breakable,colback=teal!4,colframe=teal!35!black,title={Before you move on}]
What does \fa{اتاق} mean? (\textquotedblleft{}A room\textquotedblright{}.) Say it once more.
\end{tcolorbox}
\section[\texorpdfstring{\fa{پنجره} (panjere)}{پنجره (panjere)}]{\fa{پنجره} (panjere) — a window}

\label{lesson:FA-C45-panjere}



\begin{quote}
Before the new one: say the Persian for a house, then the Persian for a room.
\end{quote}

\subsection*{You'll want to know: \fa{پنجره}}
\textbf{\fa{پنجره}} — \emph{panjere} — \textquotedblleft{}a window\textquotedblright{}.

\begin{cousinweb}
\emph{panjere} has two accounts. It is often linked with \emph{panje}, \textquotedblleft{}paw, hand\textquotedblright{}, from \emph{panj}, \textquotedblleft{}five\textquotedblright{}; EtymDB instead sets it beside Sanskrit \emph{piñjara}, \textquotedblleft{}cage\textquotedblright{}. Which is right is not settled.
\end{cousinweb}

\begin{grammarlens}[title={What you've built}]
One more word for this chapter.
\end{grammarlens}

\subsection*{Your turn}
\begin{itemize}
\raggedright
  \item \emph{Say it:} \emph{panjere}
  \item \emph{Say it:} \emph{panjere}, once more
  \item \emph{Say it:} say \emph{panjere}
  \item \emph{Recall:} say the Persian for a house, then the Persian for a room, then say \emph{panjere} again
\end{itemize}

\begin{tcolorbox}[breakable,colback=teal!4,colframe=teal!35!black,title={Before you move on}]
What does \fa{پنجره} mean? (\textquotedblleft{}A window\textquotedblright{}.) Say it once more.
\end{tcolorbox}
\section[\texorpdfstring{\fa{صندلی} (sandali)}{صندلی (sandali)}]{\fa{صندلی} (sandali) — a chair}

\label{lesson:FA-C45-sandali}



\begin{quote}
Before the new one: say the Persian for a room, then the Persian for a window.
\end{quote}

\subsection*{You'll want to know: \fa{صندلی}}
\textbf{\fa{صندلی}} — \emph{sandali} — \textquotedblleft{}a chair\textquotedblright{}.

\begin{cousinweb}
\emph{sandali} is usually traced to \emph{sandal}, \textquotedblleft{}sandalwood\textquotedblright{}, with the ending \emph{-i}; \emph{sandal} came into Persian through Arabic \emph{ṣandal}, the fragrant wood (Lane). The chair is said to be named after the wood, the account EtymDB gives.
\end{cousinweb}

\begin{grammarlens}[title={What you've built}]
One more word for this chapter.
\end{grammarlens}

\subsection*{Your turn}
\begin{itemize}
\raggedright
  \item \emph{Say it:} \emph{sandali}
  \item \emph{Say it:} \emph{sandali}, once more
  \item \emph{Say it:} say \emph{sandali}
  \item \emph{Recall:} say the Persian for a room, then the Persian for a window, then say \emph{sandali} again
\end{itemize}

\begin{tcolorbox}[breakable,colback=teal!4,colframe=teal!35!black,title={Before you move on}]
What does \fa{صندلی} mean? (\textquotedblleft{}A chair\textquotedblright{}.) Say it once more.
\end{tcolorbox}
\section[\texorpdfstring{\fa{میز} (miz)}{میز (miz)}]{\fa{میز} (miz) — a table}

\label{lesson:FA-C45-miz}



\begin{quote}
Before the new one: say the Persian for a window, then the Persian for a chair.
\end{quote}

\subsection*{You'll want to know: \fa{میز}}
\textbf{\fa{میز}} — \emph{miz} — \textquotedblleft{}a table\textquotedblright{}.

\begin{cousinweb}
\emph{miz} is a native Persian word, inherited rather than borrowed: it continues Middle Persian \emph{mēzd}, \textquotedblleft{}offering, ritual meal\textquotedblright{} (EtymDB). The word for a ritual meal is usually said to have come to mean the table it was served on.
\end{cousinweb}

\begin{grammarlens}[title={What you've built}]
One more word for this chapter.
\end{grammarlens}

\subsection*{Your turn}
\begin{itemize}
\raggedright
  \item \emph{Say it:} \emph{miz}
  \item \emph{Say it:} \emph{miz}, once more
  \item \emph{Say it:} say \emph{miz}
  \item \emph{Recall:} say the Persian for a window, then the Persian for a chair, then say \emph{miz} again
\end{itemize}

\begin{tcolorbox}[breakable,colback=teal!4,colframe=teal!35!black,title={Before you move on}]
What does \fa{میز} mean? (\textquotedblleft{}A table\textquotedblright{}.) Say it once more.
\end{tcolorbox}
